\documentclass[12pt,letterpaper]{article}

\usepackage[spanish]{babel}
\usepackage[utf8]{inputenc}
\usepackage[T1]{fontenc}
\usepackage{graphicx}
\usepackage{fancyhdr}
\usepackage{geometry}
\usepackage{amsmath}
\usepackage{amssymb}

\geometry{margin=2.5cm}
\setlength{\headheight}{52pt}


\pagestyle{fancy}
\fancyhf{}

\fancyhead[L]{
    \includegraphics[height=1.5cm]{ur_image.png}
}

\fancyhead[R]{
\begin{tabular}{r}
\textbf{Monitoria Álgebra Lineal - Vectores} \\
Gabriel Fernando Salcedo Zamora \\
\end{tabular}
}

\renewcommand{\headrulewidth}{0.4pt}

\begin{document}

\section*{Taller de vectores}

\begin{enumerate}
    %operaciones de vectores
    \item Resuelva los siguientes problemas teniendo en cuenta que $\mathbf{u} = (-1,3,-8)$, $\mathbf{v} = (2,-2,4)$ y $\mathbf{w} = (9,-5,0)$:
    \begin{enumerate}
        \item $\mathbf{u} + \mathbf{w}$
        \item $3\mathbf{v}$
        \item $2\mathbf{u} - 5\mathbf{v}$
        \item $3\mathbf{u} - 7\mathbf{v} - 4\mathbf{w}$
        \item $\alpha \mathbf{u} - \dfrac{1}{\beta} \mathbf{v}$, con $\alpha$ y $\beta$ escalares reales.
    \end{enumerate}
    
    %vectores paralelos
    \item Determine si los vectores $\mathbf{u} = (2,-1,4)$ y $\mathbf{v} = (4,-2,8)$ son paralelos. En caso afirmativo, encuentre un escalar $\alpha$ tal que $\mathbf{u}=\alpha\mathbf{v}$.
    
    %combinacion lineal
    \item Determine si el vector $(8,6,5)$ es combinación lineal de los vectores $(-1,3,2)$ y $(5,5,4)$. En caso de que sea cierto, exprésela explícitamente; en caso contrario, justifique su respuesta.
    
    %norma de un vector
    \item Calcule la norma de los siguientes vectores:
    \begin{enumerate}
        \item $\mathbf{u} = \left(\dfrac{1}{2}, -2\right)$
        \item $\mathbf{v} = (2,2,\sqrt{3})$
        \item $\mathbf{w} = (\pi,0)$
    \end{enumerate}
    
    %norma con parametro
    \item Determine $\alpha$ tal que $\|(1, \alpha, -3, 2)\| = 5$. Recuerde que $\|\mathbf{v}\| = \sqrt{v_1^2 + v_2^2 + \cdots + v_n^2}$.
    
\end{enumerate}

\newpage

\section*{Solución}

\begin{enumerate}
    % Ejercicio 1
    \item Dados $\mathbf{u} = (-1,3,-8)$, $\mathbf{v} = (2,-2,4)$ y $\mathbf{w} = (9,-5,0)$:

    \begin{enumerate}
        \item 
        \[
        \mathbf{u} + \mathbf{w} = (-1+9,\; 3-5,\; -8+0) = (8,-2,-8).
        \]

        \item 
        \[
        3\mathbf{v} = 3(2,-2,4) = (6,-6,12).
        \]

        \item 
        \[
        2\mathbf{u} - 5\mathbf{v}
        = 2(-1,3,-8) - 5(2,-2,4)
        = (-2,6,-16) - (10,-10,20)
        = (-12,16,-36).
        \]
        
        \item 
        \[
        3\mathbf{u} - 7\mathbf{v} - 4\mathbf{w}
        = (-3,9,-24) - (14,-14,28) - (36,-20,0)
        = (-53,43,-52).
        \]

        \item 
        \[
        \alpha \mathbf{u} - \frac{1}{\beta}\mathbf{v}
        = \alpha(-1,3,-8) - \frac{1}{\beta}(2,-2,4)
        = \left(-\alpha-\frac{2}{\beta},\; 3\alpha+\frac{2}{\beta},\; -8\alpha-\frac{4}{\beta}\right),
        \]
        con $\alpha,\beta \in \mathbb{R}$ y $\beta \neq 0$.
    \end{enumerate}

    % Ejercicio 2
    \item Los vectores $\mathbf{u}=(2,-1,4)$ y $\mathbf{v}=(4,-2,8)$ son paralelos si existe un escalar $\alpha$ tal que $\mathbf{u}=\alpha\mathbf{v}$.

    Comparando componentes:
    \[
    \alpha = \frac{2}{4} = \frac{-1}{-2} = \frac{4}{8} = \frac{1}{2}.
    \]

    Como el mismo escalar satisface las tres componentes, los vectores son paralelos y
    \[
    \mathbf{u}=\frac{1}{2}\mathbf{v}.
    \]

    % Ejercicio 3
    \item Buscamos escalares $\alpha,\beta \in \mathbb{R}$ tales que
    \[
    \alpha(-1,3,2) + \beta(5,5,4) = (8,6,5).
    \]

    Esto genera el sistema de ecuaciones:
    \[
    \begin{cases}
    -\alpha + 5\beta = 8 \\
    3\alpha + 5\beta = 6 \\
    2\alpha + 4\beta = 5
    \end{cases}
    \]

    Restando la primera ecuación de la segunda:
    \[
    (3\alpha + 5\beta) - (-\alpha + 5\beta) = 6 - 8 \;\Rightarrow\; 4\alpha = -2 \;\Rightarrow\; \alpha = -\frac{1}{2}.
    \]

    Sustituyendo en la primera ecuación:
    \[
    -\left(-\frac{1}{2}\right) + 5\beta = 8 \;\Rightarrow\; \frac{1}{2} + 5\beta = 8 \;\Rightarrow\; \beta = \frac{3}{2}.
    \]

    Verificamos que estos valores satisfacen la tercera ecuación:
    \[
    2\left(-\frac{1}{2}\right) + 4\left(\frac{3}{2}\right) = -1 + 6 = 5. \checkmark
    \]

    Como el sistema es consistente para las tres ecuaciones, el vector $(8,6,5)$ \textbf{sí} es combinación lineal de $(-1,3,2)$ y $(5,5,4)$:
    \[
    (8,6,5) = -\frac{1}{2}(-1,3,2) + \frac{3}{2}(5,5,4).
    \]

    % Ejercicio 4
    \item 
    \begin{enumerate}
        \item 
        \[
        \|\mathbf{u}\| = \sqrt{\left(\frac{1}{2}\right)^2 + (-2)^2}
        = \sqrt{\frac{1}{4}+4}
        = \frac{\sqrt{17}}{2}.
        \]

        \item 
        \[
        \|\mathbf{v}\| = \sqrt{2^2+2^2+(\sqrt3)^2}
        = \sqrt{4+4+3}
        = \sqrt{11}.
        \]

        \item 
        \[
        \|\mathbf{w}\| = \sqrt{\pi^2 + 0^2} = |\pi| = \pi.
        \]
    \end{enumerate}

    % Ejercicio 5
    \item 
    \[
    \|(1,\alpha,-3,2)\| = \sqrt{1^2+\alpha^2+(-3)^2+2^2}
    = \sqrt{\alpha^2+14}.
    \]

    Como $\|(1,\alpha,-3,2)\|=5$, se tiene:
    \[
    \sqrt{\alpha^2+14}=5 \;\Rightarrow\; \alpha^2+14=25 \;\Rightarrow\; \alpha^2=11.
    \]

    Por lo tanto,
    \[
    \alpha = \pm \sqrt{11}.
    \]
\end{enumerate}

\end{document}
